\documentclass[a4paper,11pt,reqno]{amsart}

\usepackage[margin=1in]{geometry}
\usepackage[T1]{fontenc}
\usepackage{newtxtext,newtxmath}
\usepackage{microtype}
\usepackage{amsmath,mathtools}
\usepackage[numbers,sort&compress]{natbib}
\usepackage{enumitem}
\usepackage{needspace}
\usepackage{etoolbox}
\usepackage{xurl}
\usepackage[hidelinks,pdfencoding=auto,psdextra]{hyperref}
\usepackage{bookmark}

\setcounter{secnumdepth}{1}
\setcounter{tocdepth}{1}
\setlength{\parindent}{1.5em}
\setlength{\parskip}{0pt}
\setlength{\emergencystretch}{3em}
\setlist[itemize]{leftmargin=2em,itemsep=0.12em,topsep=0.35em}
\setlist[enumerate]{leftmargin=2.2em,itemsep=0.12em,topsep=0.35em}
\allowdisplaybreaks[2]
\raggedbottom
\Urlmuskip=0mu plus 1mu
\frenchspacing
\tolerance=1800
\pretolerance=500
\ifdefined\pdfgentounicode
  \pdfgentounicode=1
\fi

\providecommand{\tightlist}{\setlength{\itemsep}{0pt}\setlength{\parskip}{0pt}}
\providecommand{\passthrough}[1]{#1}
\newcommand{\proofpart}[1]{%
  \par\smallskip\noindent\textit{#1.}\enspace\ignorespaces}

\theoremstyle{plain}
\newtheorem*{theoremA}{Theorem A}
\newtheorem{theorem}{Theorem}[section]
\newtheorem{lemma}[theorem]{Lemma}
\newtheorem{proposition}[theorem]{Proposition}
\newtheorem{corollary}[theorem]{Corollary}
\newtheorem*{claim}{Claim}
\theoremstyle{definition}
\newtheorem{definition}[theorem]{Definition}
\theoremstyle{remark}
\newtheorem{remark}[theorem]{Remark}

\BeforeBeginEnvironment{theoremA}{\Needspace{8\baselineskip}}
\BeforeBeginEnvironment{theorem}{\Needspace{6\baselineskip}}
\BeforeBeginEnvironment{lemma}{\Needspace{6\baselineskip}}
\BeforeBeginEnvironment{proposition}{\Needspace{6\baselineskip}}
\BeforeBeginEnvironment{corollary}{\Needspace{5\baselineskip}}
\BeforeBeginEnvironment{claim}{\Needspace{6\baselineskip}}
\BeforeBeginEnvironment{definition}{\Needspace{5\baselineskip}}
\BeforeBeginEnvironment{remark}{\Needspace{4\baselineskip}}
\BeforeBeginEnvironment{proof}{\Needspace{4\baselineskip}}
\hypersetup{
  pdftitle={Obligatory Triple Systems: An Alternative Proof},
  pdfauthor={Samuil Petkov},
  pdfsubject={An alternative proof, finite parameter consequences, and Lean verification for obligatory triple systems},
  pdfkeywords={obligatory triple system, hypergraph colouring, Levi graph, Berge cycle, uncountable chromatic number, Erdos Problem 593},
  pdfcreator={LaTeX with amsart},
  pdfproducer={pdfTeX},
  pdfdisplaydoctitle=true
}

\title[Obligatory triple systems]{Obligatory Triple Systems: An Alternative Proof}
\author{Samuil Petkov}
\date{24 July 2026}
\subjclass[2020]{Primary 05C65; Secondary 05C15, 05C63, 03E05}
\keywords{obligatory triple system, hypergraph colouring, Levi graph, Berge cycle, uncountable chromatic number, Erdős Problem 593}

\begin{document}

\begin{abstract}
We present an alternative proof of the classification of finite obligatory
triple systems, together with a Lean formalisation of the classification and
sharp finite parameter consequences.  After isolated vertices are removed, obligatoriness is
equivalent to linearity, incidence of every hyperedge-node of the Levi graph
with a bridge, and evenness of every Berge cycle; equivalently, the systems
are generated from private-vertex expansions of finite bipartite graphs and
finite edgeless systems by disjoint unions and one-point amalgamations.  The
positive direction uses a probabilistic rainbow lemma and rooted abundance,
while the intrinsic decomposition deletes selected Levi-graph bridges and
reassembles the resulting expansion pieces along a quotient forest.  The
negative direction uses a one-apex sequence lift in a base-fibre and
support-incidence form.  As a new finite consequence, a reduced obligatory
system with $m\ge1$ hyperedges and $1\le c\le m$ connected components has
order $n$ precisely when
\[
m+2(c-1)+\left\lceil2\sqrt{m-c+1}\right\rceil
\le n\le 2m+c.
\]
All arguments are in ZFC.
\end{abstract}

\maketitle

\section*{Introduction}\label{introduction}

The hosts in this problem are infinite, but the answer is decided by a small
piece of finite geometry.  A finite triple system is \emph{obligatory} if it
occurs in every triple system of uncountable chromatic number.  Erdős asked for a characterization
of these finite systems \citep{erdos1995}; the question is catalogued as
Erdős Problem~593 \citep{bloom593}.  The answer has both a constructive form
and an intrinsic form.

For a finite graph $J$, write $J^+$ for its private-vertex expansion.  Let
$\mathcal B$ be the smallest class containing $J^+$ for every finite
bipartite graph $J$ and every finite edgeless triple system, and closed under
finite disjoint unions and one-point amalgamations.  If $F$ is a triple
system, let $F^\circ$ denote the system obtained by deleting its isolated
vertices.

\begin{theoremA}\label{theorem-a-exact-characterisation}
For every finite triple system $F$, the following are equivalent:
\begin{enumerate}
\item $F$ is obligatory;
\item $F\in\mathcal B$;
\item $F^\circ$ is linear, every hyperedge-node of $I(F^\circ)$ is incident
with a bridge, and every Berge cycle of $F^\circ$ has even length.
\end{enumerate}
\end{theoremA}

Li's preprint, posted on 23 June 2026, contains the first publicly posted
complete mathematical proof of Theorem~A \citep[Theorem~1.1]{li2026}.  It
also introduced the complete-rank one-apex lift and bridge-trace method used
in the negative direction \citep[Sections~3--4 and Theorem~4.6]{li2026}.
The present paper gives a different implementation, direct positive
arguments, sharp finite parameter consequences, and a Lean formalisation of
this implementation.

The graph analogue was proved by Erdős and Hajnal: the obligatory finite
graphs are exactly the bipartite graphs
\citep[Corollary~5.6 and Theorem~7.4]{erdoshajnal1966}.  The classical
nonlinearity obstruction for uniform hypergraphs is due to Erdős, Hajnal,
and Rothschild \citep[Theorem~2, p.~532]{ehr1973}.  Obligatory triple
systems were studied further by \citet{egh1975}, \citet{komjath2001},
\citet{hajnalkomjath2008}, and \citet{komjath2008}; the expansion theorem of
\citet[Theorem~1.2]{reiher2024} supplies the positive atoms used below.

The proof is best read as two stories.  On the positive side, a
probabilistic rainbow lemma provides the local injectivity needed to force
complete bipartite expansions.  A rooted-abundance lemma then proves closure
under one-point amalgamation, including at singular uncountable cardinals.
Deleting the bridges of the Levi graph identifies each remaining active
component with the expansion of a finite bipartite graph; the quotient graph
is a forest, and its running-intersection property gives the required
amalgamation order.

On the negative side, we use the one-apex sequence-lift strategy of
\citet[Sections~3--4]{li2026}.  The lift transfers uncountable chromatic
number from a graph to a triple system while controlling every finite linear
trace.  We prove the required trace statement through expansion fibres and a
support-incidence forest, so that the fibres are joined only at cut points.
This prevents a bridge-free hyperedge-node or a Berge cycle from being
assembled across different fibres.  There are exactly three ways to fail the
intrinsic test, and each has its own host: an Erdős--Rado linear host excludes
nonlinearity, the lift of $K_{\omega_1}$ excludes a missing bridge, and a
classical Erdős--Hajnal graph of uncountable chromatic number and prescribed
odd girth excludes an odd Berge cycle.

All embeddings are injective and non-induced.  The argument is carried out
in ZFC.  The imported results used as black boxes are the de Bruijn--Erdős
compactness theorem \citep[Theorem~1, pp.~371--373]{debruijn1951}, the
Erdős--Hajnal high-odd-girth theorem
\citep[Theorem~7.4, p.~76]{erdoshajnal1966}, and the pair case of the
Erdős--Rado partition theorem
\citep[Theorem~4(i), formula~(95), p.~471]{erdosrado1956}.

\section{Preliminaries}\label{conventions-and-definitions}

A \emph{triple system} is a simple \(3\)-uniform hypergraph \(H=(V(H),E(H))\), so \(E(H)\subseteq[V(H)]^3\). A colouring \(c:V(H)\to\lambda\) is \emph{proper} when no hyperedge is monochromatic, and \(\chi(H)\) is the least cardinal \(\lambda\) admitting such a colouring. An embedding of a triple system \(F\) in \(H\) is an injective map \(f:V(F)\to V(H)\) such that \(f[e]\in E(H)\) for every \(e\in E(F)\); embeddings are not required to be induced. A finite triple system \(F\) is \emph{obligatory} if every triple system \(H\) with \(\chi(H)>\aleph_0\) contains an embedding of \(F\). A point is \emph{isolated} if it belongs to no hyperedge, and \(F^\circ\) denotes the result of deleting all isolated points.

All ordinary graphs in this paper are simple. A triple system is \emph{linear} when any two distinct hyperedges meet in at most one point. Its \emph{Levi graph} \(I(F)\) is the bipartite graph with classes \(V(F)\) and \(E(F)\), where a point-node \(p\) is adjacent to a hyperedge-node \(e\) exactly when \(p\in e\). A \emph{bridge} is an edge of a graph whose deletion increases the number of connected components. A \emph{Berge cycle of length} \(\ell\ge2\) consists of distinct points \(p_0,\ldots,p_{\ell-1}\) and distinct hyperedges \(e_0,\ldots,e_{\ell-1}\) such that \(p_i\in e_i\cap e_{i+1}\), with indices modulo \(\ell\). Equivalently, it is a simple cycle of length \(2\ell\) in \(I(F)\).

For a finite simple graph \(J\), its \emph{private-vertex expansion} \(J^+\) has vertex set
\[
V(J)\mathbin{\dot\cup}\{p_a:a\in E(J)\}
\]
and hyperedges \(\{u,v,p_{\{u,v\}}\}\) for \(\{u,v\}\in E(J)\), where the points \(p_a\) are new and pairwise distinct. A \emph{one-point amalgamation} of vertex-disjoint triple systems \(F_0,F_1\) is obtained by choosing \(x_i\in V(F_i)\), identifying \(x_0\) with \(x_1\), and making no other identifications or new hyperedges. Finite disjoint unions and finite edgeless systems include the empty cases.

\paragraph{Finite reductions.}\label{finite-reductions}

\begin{lemma}[Isolated-vertex reduction]\label{lemma-1.1-isolated-vertex-reduction}

Let \(F\) be a finite triple system and let \(F^\circ\) be obtained by deleting all isolated vertices. Then

\[
F\text{ is obligatory}\iff F^\circ\text{ is obligatory},
\]

and

\[
F\in\mathcal B\iff F^\circ\in\mathcal B.
\]

\end{lemma}
\begin{proof}
If \(F^\circ\) is obligatory and \(H\) has uncountable chromatic number, then \(H\) has infinitely many vertices. After embedding the finite system \(F^\circ\), map the isolated vertices of \(F\) injectively to unused host vertices. Containment is non-induced, so additional host edges are irrelevant. The converse follows because \(F^\circ\) is a subhypergraph of \(F\).

For membership in \(\mathcal B\), the implication \(F^\circ\in\mathcal B\Rightarrow F\in\mathcal B\) follows by taking the disjoint union with a finite edgeless system.

For the converse, use structural induction on a construction of \(F\) in \(\mathcal B\). If \(F=J^+\), delete the isolated graph vertices of \(J\) to obtain \(J_0\); then \((J^+)^\circ=J_0^+\in\mathcal B\). If \(F\) is edgeless, its reduction is the empty edgeless system. Reduction commutes with disjoint union.

Suppose finally that \(F\) is a one-point amalgamation of \(F_0\) and
\(F_1\), identifying \(x_i\in V(F_i)\). By induction,
\(F_0^\circ,F_1^\circ\in\mathcal B\). If both selected vertices are
nonisolated in their factors, then \(F^\circ\) is the one-point amalgamation
of the two reductions at those vertices. If exactly one selected vertex is
nonisolated, the isolated side disappears at the identified point and
\(F^\circ\) is the disjoint union of the two reductions. The same description
applies when both selected vertices are isolated, because the identified
point is then deleted. Hence \(F^\circ\in\mathcal B\) in every case.
\end{proof}

Thus isolated vertices may be suppressed throughout the substantive argument and restored at the end.

\begin{lemma}[Finite deletion]\label{lemma-1.2-finite-deletion}

If \(\chi(H)>\aleph_0\) and \(S\subseteq V(H)\) is finite, then

\[
\chi(H-S)>\aleph_0.
\]

\end{lemma}
\begin{proof}
Otherwise a countable colouring of \(H-S\), together with finitely many fresh colours for \(S\), would colour \(H\) countably.
\end{proof}

\paragraph{Colouring and closure tools.}\label{colouring-and-closure-tools}

\begin{lemma}[Two elementary colouring facts]\label{lemma-1.3-two-elementary-colouring-facts}

We shall use the following two facts repeatedly.

\begin{enumerate}
\def\labelenumi{\arabic{enumi}.}
\tightlist
\item
  If \(E(H)=E_1\cup\cdots\cup E_r\) and every \((V(H),E_i)\) is countably chromatic, then \(H\) is countably chromatic.
\item
  If a graph \(G\) admits an orientation in which every vertex has outdegree at most \(d<\omega\), then \(\chi(G)\le 2d+1\).
\end{enumerate}

\end{lemma}
\begin{proof}
For the first assertion, take the product of the \(r\) countable colourings.

For the second, every finite subgraph has at most \(d|V|\) edges and therefore a vertex of degree at most \(2d\). It is \((2d+1)\)-colourable by greedy deletion. The de Bruijn--Erdős compactness theorem \citep[Theorem~1]{debruijn1951} then gives a \((2d+1)\)-colouring of all of \(G\).
\end{proof}

\begin{lemma}[Closure-chain lemma]\label{lemma-1.4-closure-chain-lemma}

Let \(\kappa\) be an uncountable cardinal, let \(1\le r<\omega\), and let

\[
\Phi:[\kappa]^r\longrightarrow[\kappa]^{<\omega}
\]

have uniformly finite values. There is an increasing continuous chain

\[
\varnothing=M_0\subseteq M_1\subseteq\cdots
\subseteq M_i\subseteq\cdots\qquad(i<\operatorname{cf}\kappa)
\]

such that

\[
\bigcup_{i<\operatorname{cf}\kappa}M_i=\kappa,
\qquad |M_i|<\kappa,
\]

and every \(M_i\) is \(\Phi\)-closed:

\[
a\in[M_i]^r\quad\Longrightarrow\quad\Phi(a)\subseteq M_i.
\]

\end{lemma}
\begin{proof}
Choose sets \(A_i\subseteq\kappa\), each of cardinality below \(\kappa\), whose union is \(\kappa\). At a successor stage close \(M_i\cup A_i\) under \(\Phi\), iterating the operation \(\omega\) times. Since \(\Phi\) is finitary and finite-valued, this does not increase an infinite cardinal below \(\kappa\). At a limit stage take the union. A union of an increasing chain of sets closed under a finitary operation remains closed; and because the stage is below \(\operatorname{cf}\kappa\), its cardinality remains below \(\kappa\).
\end{proof}

Write

\[
I_i=M_{i+1}\setminus M_i.
\]

Then \((I_i)_{i<\operatorname{cf}\kappa}\) partitions \(\kappa\), and \(M_i=\bigcup_{j<i}I_j\).

\section{A bipartite subgraph lemma}\label{the-graph-lemma-needed-for-the-expansion-atom}

The next lemma is the finite complete-bipartite consequence of
\citet[Corollary~5.6, p.~72]{erdoshajnal1966}. We include its short
closure-chain proof because the same construction is used later.

\begin{lemma}[Uncountable chromatic number forces complete bipartite graphs]\label{lemma-2.1-uncountable-chromatic-number-forces-complete-bipartite-graphs}

For every positive integer \(n\), every graph of uncountable chromatic number contains \(K_{n,n}\).

\end{lemma}
\begin{proof}
Suppose otherwise, and choose a \(K_{n,n}\)-free graph \(G\) of uncountable chromatic number whose vertex cardinality \(\kappa\) is minimal.

For \(A\in[\kappa]^n\), let

\[
N(A)=\{v\in\kappa:v\text{ is adjacent to every member of }A\}.
\]

Because \(G\) is \(K_{n,n}\)-free, \(|N(A)|<n\). Apply Lemma 1.4 to \(A\mapsto N(A)\), and let \(M_i,I_i\) be the resulting chain and layers.

By minimality of \(\kappa\), every induced graph \(G[I_i]\) is countably chromatic. Let \(G^\times\) consist of the edges joining distinct layers, and orient each such edge from the later endpoint toward the earlier endpoint. If \(v\in I_i\) had \(n\) earlier neighbours, those neighbours would form an \(n\)-set \(A\subseteq M_i\), and closure would give

\[
v\in N(A)\subseteq M_i,
\]

contrary to \(v\in I_i\). Thus the orientation of \(G^\times\) has outdegree at most \(n-1\). By Lemma 1.3, \(G^\times\) is finitely colourable.

Colour each \(G[I_i]\) with colours in \(\omega\), reusing the same palette on all layers, and combine this with a finite proper colouring of \(G^\times\). The resulting countable product colouring is proper on all of \(G\), a contradiction.
\end{proof}

\section{Bipartite expansions}\label{private-vertex-expansions-of-bipartite-graphs-are-obligatory}

\paragraph{Rainbow bipartite submatrices.}\label{rainbow-bipartite-submatrices}

\begin{lemma}[Rainbow bipartite submatrices]\label{lemma-3.1-rainbow-bipartite-submatrices}

For all positive integers \(n,t\), there exists \(q=q(n,t)\) with the following property. Suppose the edges of \(K_{q,q}\) are coloured, with an arbitrary colour set, so that at every vertex each individual colour occurs on fewer than \(t\) incident edges. Then there are \(n\)-element subsets \(X'\) and \(Y'\) of the two vertex classes such that all \(n^2\) edges between \(X'\) and \(Y'\) have distinct colours.

\end{lemma}
\begin{proof}
Choose independent uniformly random \(n\)-subsets \(X'\) and \(Y'\). For every unordered pair \(P\) of distinct same-coloured edges, let \(I_P\) be the indicator that both edges of \(P\) lie in the selected \(K_{n,n}\), and put
\[
 Z=\sum_P I_P.
\]
Thus \(Z=0\) exactly when all selected edges have distinct colours.

Let \(m_\gamma\) be the number of edges of colour \(\gamma\). The local bound gives

\[
m_\gamma\le(t-1)q.
\]

At a fixed vertex and for a fixed colour \(\gamma\), fewer than \(t\)
incident edges have colour \(\gamma\). Summing
\(\binom{d_{v,\gamma}}2\le(d_{v,\gamma})(t-1)/2\) over the \(2q\)
vertices shows that the number of same-coloured pairs sharing a vertex is
at most \((t-1)q^2\). The total number of same-coloured pairs is at most

\[
\sum_\gamma\binom{m_\gamma}{2}
\le \frac12\bigl(\max_\gamma m_\gamma\bigr)\sum_\gamma m_\gamma
\le \frac12(t-1)q^3.
\]

If the two edges share a vertex in the left class, their exact selection
probability is
\[
 \frac nq\frac{(n)_2}{(q)_2};
\]
the same formula holds when the shared vertex is on the right. If the two
edges are disjoint, the exact probability is
\[
 \left(\frac{(n)_2}{(q)_2}\right)^2.
\]
For \(q\ge2n\), these are at most \(2n^3/q^3\) and \(4n^4/q^4\),
respectively. By linearity of expectation,

\[
\mathbb{E}\!\left[Z\right]
 \le\frac{2(t-1)n^3}{q}+\frac{2(t-1)n^4}{q}.
\]

This is below \(1\) for sufficiently large \(q\). Since \(Z\) is
integer-valued and nonnegative, it cannot satisfy \(Z\ge1\) for every
choice when \(\mathbb{E}\!\left[Z\right]<1\). Hence some choice has \(Z=0\), and the
selected \(K_{n,n}\) is rainbow.
\end{proof}

\paragraph{The complete bipartite expansion atom.}\label{the-complete-bipartite-expansion-atom}
\begin{proposition}[The complete bipartite expansion atom]\label{proposition-3.2-the-complete-bipartite-expansion-atom}

For every positive integer \(n\), the triple system \(K_{n,n}^{+}\) is obligatory.

\end{proposition}
\begin{proof}
Suppose not. Choose a \(K_{n,n}^{+}\)-free triple system

\[
H=(\kappa,E)
\]

with \(\chi(H)>\aleph_0\) and with \(\kappa\) minimal.

For a pair \(xy\), write

\[
d_H(xy)=|\{z:\{x,y,z\}\in E\}|.
\]

Set \(t=3n^2+1\) and define a graph \(R\) on \(\kappa\) by

\[
xy\in E(R)\quad\Longleftrightarrow\quad d_H(xy)\ge t.
\]

\medskip\noindent\textbf{Claim 3.2.1 (the graph \(R\) is countably chromatic).}

Otherwise Lemma 2.1 gives a \(K_{n,n}\) in \(R\), with core vertex set \(C\), \(|C|=2n\). For each of its \(n^2\) graph edges \(xy\), choose a third vertex \(p_{xy}\) with

\[
\{x,y,p_{xy}\}\in E,
\]

avoiding \(C\) and all previously chosen private vertices. At every stage fewer than \(2n+n^2<t\) vertices are forbidden, while \(xy\) has at least \(t\) possible third vertices. Thus the \(p_{xy}\) can be chosen distinct and outside \(C\), producing a copy of \(K_{n,n}^{+}\), a contradiction. \(\square\)

Let

\[
E_1=\{e\in E:\text{some pair contained in }e\text{ is an edge of }R\},
\qquad E_2=E\setminus E_1.
\]

A proper colouring of \(R\) is a proper colouring of \((\kappa,E_1)\), since every \(E_1\)-edge contains an \(R\)-edge. Hence \(\chi(\kappa,E_1)\le\aleph_0\), and therefore

\[
\chi(\kappa,E_2)>\aleph_0. \tag{3.1}
\]

For a pair \(xy\), define

\[
\Phi(xy)=
\begin{cases}
\{z:\{x,y,z\}\in E\},&d_H(xy)<t,\\
\varnothing,&d_H(xy)\ge t.
\end{cases}
\]

Apply Lemma 1.4 to \(\Phi\), obtaining \(M_i,I_i\).

Every \(e\in E_2\) has at least two vertices in its highest layer. Indeed, suppose the highest layer met \(e\) only in \(z\in I_i\). Then the other pair \(xy=e\setminus\{z\}\) lies in \(M_i\). Since \(e\in E_2\), \(xy\notin R\), so \(d_H(xy)<t\). Closure gives \(z\in\Phi(xy)\subseteq M_i\), a contradiction.

Thus every \(E_2\)-edge is either contained in one layer, or has exactly two vertices in its highest layer \(I_i\) and its third vertex in \(M_i\). Let \(E_i^*\) be the \(E_2\)-edges contained in \(I_i\). Since \(|I_i|<\kappa\) and \(H\) is \(K_{n,n}^{+}\)-free, minimality of \(\kappa\) gives \(\chi(I_i,E_i^*)\le\aleph_0\). Consequently

\[
E^*=\bigcup_iE_i^*
\]

is countably chromatic. By (3.1), the remaining crossing-edge system

\[
E^{**}=E_2\setminus E^*
\]

has uncountable chromatic number.

For each \(i\), define a graph \(G_i\) on \(I_i\) by declaring \(xy\in E(G_i)\) when there is a \(\beta\in M_i\) such that \(\{x,y,\beta\}\in E^{**}\). If every \(G_i\) were countably chromatic, colour each \(I_i\) properly in \(G_i\), using the same countable palette on every layer. Every edge of \(E^{**}\) would then have differently coloured top-layer vertices, giving a countable colouring of \((\kappa,E^{**})\), a contradiction. Hence some \(G_i\) is uncountably chromatic.

Choose \(q=q(n,t)\) from Lemma 3.1. Lemma 2.1 gives a \(K_{q,q}\) in \(G_i\), with sides \(X,Y\). For every \(xy\in X\times Y\), choose

\[
\beta_{xy}\in M_i
\]

such that \(\{x,y,\beta_{xy}\}\in E^{**}\), and regard \(\beta_{xy}\) as the colour of the edge \(xy\).

This edge-colouring is locally \((t-1)\)-bounded. For example, if for fixed \(x\) and fixed \(\beta\) there were \(t\) different vertices \(y\) with \(\beta_{xy}=\beta\), then the pair \(\{x,\beta\}\) would have codegree at least \(t\). It would belong to \(R\), forcing the triples \(\{x,y,\beta\}\) into \(E_1\), contrary to their membership in \(E^{**}\subseteq E_2\). The same argument applies at vertices of \(Y\).

Lemma 3.1 gives \(n\)-sets \(X'\subseteq X\), \(Y'\subseteq Y\) for which the \(\beta_{xy}\), \(x\in X'\), \(y\in Y'\), are all distinct. They lie in \(M_i\), whereas \(X'\cup Y'\subseteq I_i\), so none is a core vertex. The triples

\[
\{x,y,\beta_{xy}\},
\qquad x\in X',\ y\in Y',
\]

form \(K_{n,n}^{+}\), the final contradiction.
\end{proof}

\begin{corollary}[Every finite bipartite expansion is obligatory]\label{corollary-3.3-every-finite-bipartite-expansion-is-obligatory}

If \(J\) is any finite bipartite graph, then \(J^+\) is obligatory.

\end{corollary}
\begin{proof}
Embed the two vertex classes of \(J\) in the two classes of some \(K_{n,n}\). Then \(J^+\) is a subhypergraph of \(K_{n,n}^+\). Obligatoriness is downward closed under non-induced containment: a copy of the larger finite system contains a copy of the smaller one. Apply Proposition 3.2.
\end{proof}

Corollary~3.3 is also a consequence of
\citet[Theorem~1.2]{reiher2024}; the preceding argument gives a direct proof
in the present notation.

\section{Closure properties}\label{closure-of-obligatoriness-under-the-operations-defining-mathcal-b}

Closure under disjoint unions and one-point amalgamations is part of the
established theory; see \citet[pp.~233--238]{komjath2001} and the summary in
\citet[p.~1]{reiher2024}. We include complete proofs because the rooted
abundance argument is used below.

\paragraph{Disjoint-union closure.}\label{disjoint-union-closure}

\begin{lemma}[Disjoint-union closure]\label{lemma-4.1-disjoint-union-closure}

The class of finite obligatory triple systems is closed under finite disjoint unions.

\end{lemma}
\begin{proof}
Let \(F_1,F_2\) be obligatory and let \(H\) be uncountably chromatic. First find a copy of \(F_1\). Delete its finite vertex set. By Lemma 1.2 the remaining triple system is still uncountably chromatic, so it contains \(F_2\). The two copies are disjoint. Iteration proves the finite case.
\end{proof}

\paragraph{Rooted abundance and one-point amalgamation.}\label{rooted-abundance-and-one-point-amalgamation}

\begin{definition}[Rooted abundance]\label{definition-4.2-rooted-abundance}

Let \((F,r)\) be a finite rooted triple system with \(|V(F)|>1\). For a triple system \(H\) and \(m\ge1\), let \(R_m(F,r;H)\) be the set of vertices \(v\) for which there are \(m\) rooted copies of \(F\), all sending \(r\) to \(v\), whose off-root vertex sets are pairwise disjoint.

\end{definition}
\begin{lemma}[Rooted-abundance lemma]\label{lemma-4.3-rooted-abundance-lemma}

If \((F,r)\) is finite, \(|V(F)|>1\), and \(F\) is obligatory, and if
\(\chi(H)>\aleph_0\), then

\[
\chi\bigl(H[R_m(F,r;H)]\bigr)>\aleph_0.
\]

\end{lemma}
\begin{proof}
Put

\[
B=V(H)\setminus R_m(F,r;H).
\]

For each \(v\in B\), consider the family of off-root vertex sets of all
copies of \(F\) rooted at \(v\). Because
\(v\notin R_m(F,r;H)\), this family has matching number at most \(m-1\).
Choose a maximal pairwise disjoint subfamily. It contains at most \(m-1\)
members; let \(S_v\) be their union. Then

\[
|S_v|\le(m-1)(|V(F)|-1)=:D.
\]

Moreover, \(S_v\) meets the off-root part of every rooted copy at \(v\):
otherwise that off-root set would be disjoint from every member of the
chosen subfamily and could be added to it, contradicting maximality.

Form a graph \(D_B\) on \(B\) by joining \(v\) to each member of \(S_v\cap B\). If an edge is generated in both directions, assign it to one of its generators, and orient it away from that generator. Every vertex has outdegree at most \(D\), so Lemma 1.3 gives a finite proper colouring of \(D_B\).

Let \(A\) be one colour class. Then

\[
S_v\cap A=\varnothing\qquad(v\in A).
\]

Consequently \(H[A]\) contains no copy of \(F\). Indeed, if \(\varphi\) were such a copy and \(v=\varphi(r)\), its off-root part would have to meet \(S_v\) although both are subsets of \(A\). Since \(F\) is obligatory, every \(F\)-free triple system is countably chromatic. Thus each of the finitely many colour classes of \(D_B\) induces a countably chromatic subsystem. Tagging a vertex by its finite \(D_B\)-colour and its countable colour inside that class gives a countable proper colouring of \(H[B]\).

If \(H[R_m(F,r;H)]\) were countably chromatic as well, tagging the two sides would give a countable colouring of \(H\). This contradicts \(\chi(H)>\aleph_0\).
\end{proof}

\begin{proposition}[One-point amalgamation closure]\label{proposition-4.4-one-point-amalgamation-closure}

The class of finite obligatory triple systems is closed under one-point amalgamations.

\end{proposition}
\begin{proof}
Let \(F_1,F_2\) be obligatory, with selected vertices \(r_1,r_2\), and let \(F\) be their one-point amalgamation. If one factor consists only of its selected vertex, the assertion is immediate. Otherwise set

\[
m=|V(F_2)|.
\]

By Lemma 4.3, the set

\[
R=R_m(F_1,r_1;H)
\]

induces an uncountably chromatic subsystem of every uncountably chromatic \(H\). Hence \(H[R]\) contains a copy of \(F_2\); write \(\varphi\) for its embedding and let \(v=\varphi(r_2)\).

At \(v\) there are \(m\) rooted copies of \(F_1\) with pairwise disjoint off-root sets. The set \(\varphi(V(F_2)\setminus\{r_2\})\) has \(m-1\) vertices, so it meets at most \(m-1\) of those off-root sets. One rooted \(F_1\)-copy therefore avoids the \(F_2\)-copy outside \(v\). Their union is the required one-point amalgamation.
\end{proof}

Corollary 3.3, Lemma 4.1, Proposition 4.4, and the trivial edgeless case prove

\[
(2)\Longrightarrow(1). \tag{4.1}
\]

\section{The intrinsic decomposition}\label{equivalence-of-the-constructive-and-intrinsic-descriptions}

We give the selected-incidence decomposition and running-intersection
argument in the explicit form used by the Lean development; compare the
structural framework in \citet[Sections~4--5]{li2026}.

We prove \((2)\Longleftrightarrow(3)\). Isolated vertices may be ignored by Lemma 1.1.

\paragraph{Forward preservation.}\label{forward-preservation}

\begin{proposition}[The intrinsic conditions are preserved by the generators and operations]\label{proposition-5.1-the-intrinsic-conditions-are-preserved-by-the-generators-and-operations}

Every generator of \(\mathcal B\) satisfies the intrinsic conditions, and the conditions are preserved under finite disjoint unions and one-point amalgamations.

\end{proposition}
\begin{proof}
For a bipartite graph \(J\), the expansion \(J^+\) is linear. Every hyperedge-node is joined to its private point by a bridge. A Berge cycle cannot use a private point, so the Berge cycles of \(J^+\) are precisely the ordinary graph cycles of \(J\), with the same lengths; these are even. Edgeless systems satisfy the conditions vacuously, and disjoint union plainly preserves them.

Consider a one-point amalgamation. Hyperedges belonging to different factors intersect only in the identified point, so linearity is preserved. Its Levi graph is the vertex-sum of the two Levi graphs at a point-node. Every simple cycle in a vertex-sum lies wholly in one factor: a cycle crossing into the other factor would have to pass through the cut vertex twice. Consequently bridges in the factors remain bridges, and every Berge cycle lies in one factor and retains its parity.
\end{proof}

Thus

\[
(2)\Longrightarrow(3). \tag{5.1}
\]

\paragraph{Bridge-block reconstruction.}\label{bridge-block-reconstruction}

\begin{proposition}[Bridge-block decomposition]\label{proposition-5.2-bridge-block-decomposition}

Suppose \(F\) is finite, has no isolated vertices, and satisfies the three intrinsic conditions. Then \(F\in\mathcal B\).

\end{proposition}
\begin{proof}
If \(E(F)=\varnothing\), then the assumption that \(F\) has no isolated vertices gives \(V(F)=\varnothing\), and \(F\) is an edgeless generator. Hence assume \(E(F)\ne\varnothing\). Put \(L=I(F)\), let \(B(L)\) be its set of bridges, and let \(\mathscr C\) be the set of connected components of \(L-B(L)\).

Consider a hyperedge-node \(e\). Since \(e\) is incident with a bridge and has degree three, it has at most two nonbridge incidences. It cannot have exactly one: a nonbridge incidence lies on a cycle, and that cycle must leave \(e\) through a second nonbridge incidence. Therefore every hyperedge-node has either no nonbridge incidences or exactly two.

\paragraph{The nontrivial components.}\label{the-nontrivial-components}

Let \(C\in\mathscr C\) contain a nonbridge incidence. Every hyperedge-node \(e\in C\) has exactly two point-neighbours in \(C\). Contract each such \(e\) to an ordinary graph edge between those two point-neighbours. This produces a finite graph \(J_C\).

The graph \(J_C\) is simple: parallel edges would correspond to two hyperedges of \(F\) sharing two points, contrary to linearity. Every cycle in \(J_C\) gives a Berge cycle of the same length in \(F\). Hence \(J_C\) has no odd cycle and is bipartite.

Each hyperedge-node \(e\in C\) has a third point-neighbour \(p_e\) across a bridge. These third points satisfy:

\begin{enumerate}
\def\labelenumi{\arabic{enumi}.}
\tightlist
\item
  \(p_e\notin C\); otherwise a path in \(C\) from \(p_e\) to \(e\), together with the incidence \(ep_e\), would place that incidence on a cycle.
\item
  If \(e\ne f\) lie in \(C\), then \(p_e\ne p_f\); otherwise a path in \(C\) from \(e\) to \(f\), together with \(ep_e\) and \(p_ff\), would place those bridge incidences on a cycle.
\end{enumerate}

Define a map from \(J_C^+\) to the subsystem \(P_C\) formed by these
hyperedges as follows. Send every graph vertex of \(J_C\) to the
corresponding point-node of \(C\), send the private vertex indexed by the
contracted edge associated with \(e\) to \(p_e\), and send that expanded edge
to the hyperedge \(e\) of \(F\). The two observations above show that the
vertex map is injective: the private points are pairwise distinct and lie
outside \(C\). Every \(e\) has exactly the two contracted endpoints and
\(p_e\) as its three vertices, so this map is incidence-preserving and
surjective on the vertices and hyperedges of \(P_C\). Hence it is an
isomorphism \(J_C^+\cong P_C\).

\paragraph{Components consisting of one hyperedge-node.}\label{components-consisting-of-one-hyperedge-node}

If a hyperedge-node \(e\) has no nonbridge incidence, it is an isolated component after the bridges are deleted. Its single triple is the expansion of a one-edge bipartite graph. Thus every hyperedge of \(F\) belongs to a piece \(P_C\) isomorphic to \(J^+\) for a finite bipartite graph \(J\). The hyperedge sets of these pieces partition \(E(F)\). Their vertex sets cover \(V(F)\): every point is incident with a hyperedge, and it either lies in that hyperedge-node's bridge-deleted component or is the point endpoint of a bridge leaving it.

\paragraph{Quotient forest.}\label{quotient-forest}

Contract every component of \(L-B(L)\). The original bridges become edges of a quotient multigraph \(T\). Two distinct bridges cannot join the same pair of bridge-deleted components: paths inside those two components, together with the two bridges, would form a cycle containing both, contradicting that they are bridges. Thus \(T\) is a simple graph.

To see that \(T\) is a forest, suppose that it contained a cycle. Within
each contracted component on this cycle, choose a path joining the
endpoints of the two incident quotient edges. Concatenating these paths
with the corresponding bridges gives a closed walk in \(L\) that traverses
each of those bridges exactly once. This is impossible: deleting a bridge
separates its endpoints, so every closed walk crosses the resulting cut an
even number of times. Hence \(T\) is a forest.

A vertex \(C\in V(T)=\mathscr C\) is called \emph{active} if the bridge-deleted component \(C\) contains a hyperedge-node. For an active \(C\), the piece \(P_C\) consists of the hyperedges represented in \(C\), all point-nodes of \(C\), and the point endpoints of bridge incidences leaving hyperedge-nodes in \(C\). An inactive component contains no hyperedge-node and hence no incidence of \(L-B(L)\), so it is a singleton point-node.

At this stage every hyperedge already belongs to a bipartite expansion piece.
The only remaining issue is global: the pieces must be ordered so that a new
piece never meets the assembled system in two different points. The quotient
forest records exactly the attachment geometry needed for this running
intersection.

\begin{claim}[5.2.1. Running intersection]

Every component of \(T\) contains an active vertex, because \(F\) has no isolated vertices. Choose an active root in each component, and order the active vertices by nondecreasing depth, with arbitrary order among equal-depth vertices. When an active component \(C\) is added, its piece \(P_C\) meets the union of the earlier pieces in at most one point.

\end{claim}
\begin{proof}[Proof of the claim]
For a point \(p\in V(F)\), let \(X(p)\in\mathscr C\) be the component of \(L-B(L)\) containing the point-node \(p\). The active pieces containing \(p\) are precisely \(P_{X(p)}\), when \(X(p)\) is active, together with the pieces \(P_D\) at active neighbours \(D\) of \(X(p)\) whose corresponding original bridge has point endpoint \(p\). Thus the vertices of \(T\) indexing pieces that contain \(p\) are contained in the closed star centred at \(X(p)\).

Suppose that \(p\in P_C\) also lies in an earlier piece. If \(X(p)=C\),
every other piece containing \(p\) is indexed by a neighbour of \(C\). A tree
has only one neighbour of \(C\) at smaller depth, namely its parent, so only
that piece can be earlier. If \(X(p)\ne C\), then \(C\) is a neighbour of
\(X(p)\). The vertex \(X(p)\) cannot be a child of \(C\), because then both
\(X(p)\) and every other vertex in its closed star would be deeper than
\(C\). Hence \(X(p)\) is the parent of \(C\). In either case, \(p\) is the
point endpoint of the unique parent edge of \(C\). The same reasoning applies
to every point shared with an earlier piece, so all such points coincide with
this parent point \(p_C\), and

\[
P_C\cap\bigcup_{D\text{ earlier}}P_D\subseteq\{p_C\}.
\]

For a root active component the intersection is empty.
\end{proof}

Starting with the root piece in each component, add the active pieces in this
order. If the new intersection is empty, use a disjoint union. If it is the
singleton \(\{p_C\}\), use a one-point amalgamation at that point. Several
earlier pieces may already contain the same inactive attachment point, but
the new step still identifies only that one existing vertex.

For completeness, the identity maps from the pieces into \(F\) assemble
inductively to an incidence isomorphism from the constructed system onto the
union of the pieces already added. Indeed, the running-intersection claim
shows that the next piece and the preceding union have either no common
vertex or exactly the one vertex used in the amalgamation; hence the
construction makes no unintended identification. The piece hyperedge sets
are disjoint and partition \(E(F)\), while their vertex sets cover \(V(F)\).
After all active pieces have been added, the resulting subsystem is therefore
all of \(F\). Different components of \(T\) are combined by disjoint union.
Hence \(F\in\mathcal B\).
\end{proof}

This proves

\[
(3)\Longrightarrow(2). \tag{5.2}
\]

Combining (5.1) and (5.2),

\[
(2)\Longleftrightarrow(3). \tag{5.3}
\]

\section{The sequence lift}\label{the-sequence-lift-used-in-the-avoidance-constructions}

We use the one-apex construction and bridge-trace strategy of
\citet[Sections~3--4 and Theorem~4.6]{li2026}, specialised to $\omega_1$ and
written as a base-fibre and support-incidence decomposition convenient for
formalisation.

\paragraph{The one-apex construction.}\label{the-one-apex-construction}

Let \(G=(X,A)\) be an arbitrary graph, where \(A=E(G)\) is also regarded as an alphabet. Let

\[
T=A^{<\omega_1}
\]

be the tree of all sequences of graph edges of countable ordinal length. For \(s\in T\) and \(x\in X\), write \(x_s=(s,x)\).

\begin{definition}[The one-apex lift]\label{definition-6.1-the-one-apex-lift}

Define the triple system \(\mathcal L(G)\) by

\[
V(\mathcal L(G))=T\times X.
\]

For \(s\in T\), an edge \(a=\{x,y\}\in A\), a sequence \(t\in T\) extending \(s^\frown a\), and \(z\in X\), put

\[
\{x_s,y_s,z_t\}\in E(\mathcal L(G)). \tag{6.1}
\]

The two vertices \(x_s,y_s\) form the base pair and \(z_t\) is the apex. Since \(t\) properly extends \(s\), these are three distinct vertices. In the resulting triple, the sequence node occurring twice uniquely identifies \(s\); its two graph labels identify the edge \(a=\{x,y\}\), and the remaining vertex identifies \(t,z\). Thus each hyperedge arises from unique data, up to interchanging \(x\) and \(y\), and the system is simple and \(3\)-uniform.

\end{definition}

\paragraph{Chromatic and finite-trace properties.}\label{chromatic-and-finite-trace-properties}

\begin{lemma}[Chromatic lower bound]\label{lemma-6.2-chromatic-lower-bound}

If \(\chi(G)>\aleph_0\), then

\[
\chi(\mathcal L(G))>\aleph_0.
\]

\end{lemma}
\begin{proof}
Suppose \(c:T\times X\to\omega\) is proper. For every \(s\in T\), the restriction \(x\mapsto c(x_s)\) is a countable colouring of \(G\). Hence there is an edge

\[
a_s=\{x_s^0,x_s^1\}\in E(G)
\]

whose two vertices at node \(s\) have a common colour \(k_s\).

Construct \(s_\alpha\in A^\alpha\) for \(\alpha<\omega_1\) by

\[
s_0=\varnothing,
\qquad s_{\alpha+1}=s_\alpha^\frown a_{s_\alpha},
\]

and at limit ordinals take unions.

If \(\alpha<\beta\), then \(s_\beta\) extends \(s_\alpha^\frown a_{s_\alpha}\). If \(k_{s_\alpha}=k_{s_\beta}\), take either endpoint \(z\) of the monochromatic edge at \(s_\beta\). By (6.1),

\[
\{(s_\alpha,x_{s_\alpha}^0),
  (s_\alpha,x_{s_\alpha}^1),
  (s_\beta,z)\}
\]

is a monochromatic hyperedge. Therefore the colours \(k_{s_\alpha}\), \(\alpha<\omega_1\), are pairwise distinct, impossible in \(\omega\).
\end{proof}

The next result is the fibre-decomposition form of Li's bridge-trace theorem
needed in this paper \citep[Theorem~4.6]{li2026}.

\begin{theorem}[Finite linear trace decomposition]\label{theorem-6.3-exact-finite-linear-trace-theorem}

Every finite, not necessarily induced, linear subhypergraph \(K\) of
\(\mathcal L(G)\) is obtainable, by finite disjoint unions and one-point
amalgamations, from systems \(J^+\) where \(J\) is a finite subgraph of
\(G\).

\end{theorem}
\begin{proof}
The following dictionary records the three local objects used
in the assembly. A \emph{base node} is a sequence \(s\) at which a lift
hyperedge begins. Its two vertices with sequence coordinate \(s\) are the
\emph{core vertices}; its third vertex, whose coordinate properly extends
\(s\), is the \emph{apex}. The \emph{base fibre} \(K_s\) consists of all
hyperedges of \(K\) based at \(s\), together with their incident vertices.
An \emph{active base node} is one whose fibre is nonempty. The proof first
identifies each fibre with one expansion atom and then shows that the fibres
are attached along a forest of one-point overlaps.

First set aside every isolated vertex of \(K\). Each is a one-vertex edgeless
factor, hence isomorphic to \(J^+\) for a one-vertex edgeless subgraph \(J\)
of \(G\), and may be restored by disjoint union at the end. We may therefore
assume that every vertex of \(K\) lies in a hyperedge.

Every hyperedge of \(\mathcal L(G)\) has a unique base node: exactly two of
its vertices have sequence coordinate \(s\), while the apex has a strictly
longer sequence coordinate. We group the hyperedges of \(K\) by this
coordinate and use the notation \(K_s\) from the dictionary.

\medskip
\noindent\textbf{Claim 6.3.1 (each base fibre is one expansion atom).}
For every active base node \(s\), there is a finite subgraph \(J_s\) of
\(G\) such that \(K_s\cong J_s^+\).

\smallskip
\begin{proof}[Proof of Claim 6.3.1]

For a fixed graph edge \(a=\{x,y\}\in E(G)\), at most one hyperedge of \(K_s\) can have base pair \(\{x_s,y_s\}\), since two such hyperedges would share two vertices and violate linearity. Let \(J_s\) be the finite subgraph of \(G\) whose edges are precisely the graph edges used by hyperedges of \(K_s\), with vertex set consisting of their endpoints.

Define a map \(J_s^+\to K_s\) by sending each graph vertex \(x\in V(J_s)\)
to the core vertex \(x_s\), and by sending the private vertex of
\(a=\{x,y\}\) to the apex of the unique hyperedge of \(K_s\) with base pair
\(\{x_s,y_s\}\). Distinct graph edges have distinct apices: if \(z_t\)
were the apex for \(a\) and \(a'\), the coordinate of \(t\) immediately
after \(s\) would have to equal both \(a\) and \(a'\). No apex is a core
vertex, because an apex has sequence coordinate strictly extending \(s\),
whereas every core vertex has coordinate \(s\). The map is therefore
injective, and by construction it is bijective on hyperedges and preserves
each incidence. Hence \(K_s\cong J_s^+\).
\end{proof}

\medskip
\noindent\textbf{Claim 6.3.2 (two fibres have at most one common point).}
For distinct active base nodes \(s\) and \(u\),
\[
|V(K_s)\cap V(K_u)|\le 1.
\]

\smallskip
\begin{proof}[Proof of Claim 6.3.2]
If a point \(z_t\) lies in both \(K_s\) and \(K_u\), then both \(s\) and
\(u\) are prefixes of \(t\), so they are comparable. Suppose
\(s\subsetneq u\). A common point cannot be a core point of \(K_s\); it is
an apex of \(K_s\). The graph edge used by its hyperedge in \(K_s\) is the
coordinate of \(u\) immediately following \(s\). Hence all points in
\(K_s\cap K_u\) would be apices of the same base edge of \(K_s\). There is
at most one such apex. Therefore

\[
|V(K_s)\cap V(K_u)|\le1. \tag{6.2}
\]
\end{proof}

Let \(Q\) be the bipartite incidence graph of the base fibres. Its left
class is the set of active base nodes. Its right class consists of points
that lie in at least two fibres. Join a base node \(s\) to a point \(p\)
exactly when \(p\in V(K_s)\).

\medskip
\noindent\textbf{Claim 6.3.3 (the overlap graph is a forest).}
The graph \(Q\) is acyclic.

\smallskip
\begin{proof}[Proof of Claim 6.3.3]
Suppose that \(Q\) has a cycle, and choose on it a base node \(s\) of
minimum sequence length. Its two neighbouring base nodes on the cycle are
proper descendants of \(s\). The two intervening shared points are distinct,
so they are distinct apices in \(K_s\), arising from two distinct graph edges
\(a\ne a'\). Thus the two neighbouring base nodes lie in different immediate
branches below \(s\), one beginning with \(a\) and the other with \(a'\).

Starting from one neighbour of \(s\), every base node on the remaining
part of the cycle lies in the same immediate branch below \(s\). Indeed,
consecutive base nodes share a point and are therefore comparable. For two
comparable proper descendants of \(s\), the first sequence coordinate after
\(s\) is the same, so a length-two step through a shared point cannot move
from one immediate branch to another. Induction along the remaining path
therefore keeps every base node in the first branch, contradicting that its
final node lies in the other branch. Hence \(Q\) is acyclic.
\end{proof}

Root each component of \(Q\) at a base node, and order its base nodes by
nondecreasing even distance from the root. For a nonroot base node \(s\), let
\(p_s\) be the point-node adjacent to \(s\) on the path to the root.

\begin{claim}[6.3.4. Base-fibre running intersection]

Every earlier factor that meets \(K_s\) meets it at \(p_s\).

\end{claim}
\begin{proof}[Proof of Claim 6.3.4]
Suppose \(K_u\) is earlier and contains a point
\(p\in K_s\). Then \(Q\) contains the length-two path \(s-p-u\). Let the
distance of \(s\) from the root be \(2d\). The point-node \(p\) has distance
either \(2d-1\) or \(2d+1\). In the second case \(u\) has distance \(2d+2\),
so it cannot be earlier. Thus \(p\) has distance \(2d-1\), and the uniqueness
of the path to the root in the tree \(Q\) gives \(p=p_s\). This also covers
an earlier sibling \(u\), which shares the same parent point \(p_s\).
\end{proof}

\medskip
\noindent\textbf{Claim 6.3.5 (the ordered fibres assemble to \(K\)).}
Let \(U_s\) be the union of the factors preceding \(K_s\). Equation (6.2)
and Claim 6.3.4 give
\[
V(K_s)\cap U_s\subseteq\{p_s\}.
\]
Starting with the root factor, add the remaining factors in this order by
disjoint unions or one-point amalgamations as appropriate. At each step the
natural maps into \(K\) identify exactly the displayed common point and
nothing else, so they assemble to an isomorphism onto the union of the
factors already added. Different components of \(Q\), and the isolated
factors removed at the start, are combined by disjoint union. This is the
running-intersection step: Claim 6.3.1 supplies the expansion atoms, while
Claims 6.3.2--6.3.4 certify that each new atom is attached along at most one
previous point.
\end{proof}

\begin{corollary}[Restrictions on finite linear traces]\label{corollary-6.4-restrictions-on-finite-linear-traces}

Every finite linear \(K\subseteq\mathcal L(G)\) satisfies:

\begin{enumerate}
\def\labelenumi{\arabic{enumi}.}
\tightlist
\item
  every hyperedge-node of \(I(K)\) is incident with a bridge;
\item
  every Berge cycle of \(K\) is contained in a single factor \(J_s^+\);
\item
  a Berge cycle of length \(\ell\) in \(K\) yields an ordinary cycle of length \(\ell\) in \(G\).
\end{enumerate}

\end{corollary}
\begin{proof}
In every \(J_s^+\) the incidence with the private vertex is a bridge. Bridgehood is preserved by one-point amalgamation. A simple Levi cycle cannot cross a one-point cut, so every Berge cycle lies in one factor. Berge cycles in \(J_s^+\) are exactly ordinary graph cycles in \(J_s\).
\end{proof}

\section{A classical high-odd-girth input}\label{uncountably-chromatic-graphs-of-large-odd-girth}

For the odd-cycle obstruction we use the following older graph theorem
directly, rather than a shift-graph surrogate.

\begin{theorem}[Erdős--Hajnal]\label{theorem-7.1-high-odd-girth}
For every positive integer $m$ there is a graph $G$ such that
\[
\chi(G)>\aleph_0
\]
and $G$ contains no odd cycle of length at most $m$.
\end{theorem}

\begin{proof}[Source]
Theorem~7.4 of \citet[p.~76]{erdoshajnal1966}, as quoted in
\citet[Theorem~C, p.~428]{egh1975}, gives for every positive integer $i$ an
uncountably chromatic graph containing no cycle $C_{2j+1}$ for $0<j<i$.
Choose $i$ so that $2i-1\ge m$.
\end{proof}

\section{Avoidance hosts}\label{the-avoidance-hosts}

The division into nonlinear, missing-bridge, and odd-cycle hosts follows the
negative-half architecture of \citet[Sections~2--4]{li2026}.  The first two
hosts are written out in the present notation; for the third we invoke the
older Erdős--Hajnal theorem from Section~7 directly.

We prove

\[
\neg(3)\Longrightarrow\neg(1).
\]

\paragraph{Nonlinearity obstruction.}\label{nonlinearity-obstruction}

\begin{proposition}[Avoidance of every nonlinear finite triple system]\label{proposition-8.1-avoidance-of-every-nonlinear-finite-triple-system}

There is an uncountably chromatic linear triple system. Consequently, every finite nonlinear triple system is non-obligatory.

\end{proposition}
\begin{proof}
We use the Erd\H{o}s--Rado relation
\citep[Theorem~4(i), formula~(95), p.~471]{erdosrado1956}

\[
(2^{\aleph_0})^+\longrightarrow(\aleph_1)^2_{\aleph_0}. \tag{8.1}
\]

Let \(\kappa=(2^{\aleph_0})^+\). Define a triple system \(T_\kappa\) whose vertices are the pairs in \([\kappa]^2\), and whose hyperedges are

\[
\bigl\{
\{\alpha,\beta\},
\{\alpha,\gamma\},
\{\beta,\gamma\}
\bigr\},
\qquad \alpha<\beta<\gamma<\kappa.
\]

This triple system is linear: two distinct triangles of a complete graph cannot share two graph edges, because two graph edges of a triangle determine the third edge and the three underlying vertices.

It is uncountably chromatic. A countable vertex-colouring of \(T_\kappa\) is a countable edge-colouring of \(K_\kappa\). By (8.1) there is a monochromatic subset of cardinality \(\aleph_1\), and in particular a monochromatic graph triangle. This is a monochromatic hyperedge of \(T_\kappa\).

If a finite triple system \(F\) is nonlinear, it has two distinct edges sharing at least two vertices. No injective embedding of those edges can exist in the linear host \(T_\kappa\). Thus \(T_\kappa\) is \(F\)-free.
\end{proof}

\paragraph{Missing-bridge obstruction.}\label{missing-bridge-obstruction}

\begin{proposition}[Avoidance when the bridge condition fails]\label{proposition-8.2-avoidance-when-the-bridge-condition-fails}

Let \(F\) be finite and linear, and suppose some hyperedge-node of \(I(F^\circ)\) is incident with no bridge. Then \(F\) is non-obligatory.

\end{proposition}
\begin{proof}
This is the $K_{\omega_1}$ specialisation of Li's sequence-lift
missing-bridge obstruction \citep[Section~4]{li2026}.  By Lemma 1.1 it is
enough to avoid \(F^\circ\), so assume that \(F\) has no isolated vertices.
Take \(G=K_{\omega_1}\). Then \(\chi(G)>\aleph_0\), and Lemma 6.2 gives

\[
\chi(\mathcal L(G))>\aleph_0.
\]

If \(F\) appeared in \(\mathcal L(G)\), retain exactly the host hyperedges corresponding to the edges of \(F\). Because the embedding is injective and \(F\) is linear, these selected host hyperedges form a finite linear subhypergraph isomorphic to \(F\). Corollary 6.4 says that every hyperedge-node of such a trace is incident with a bridge, contradicting the chosen hyperedge of \(F\). Therefore \(\mathcal L(K_{\omega_1})\) is an uncountably chromatic \(F\)-free host.
\end{proof}

\paragraph{Odd-cycle obstruction.}\label{odd-cycle-obstruction}

\begin{proposition}[Avoidance of an odd Berge cycle]\label{proposition-8.3-avoidance-of-an-odd-berge-cycle}

Let \(F\) be finite and linear, and suppose \(F^\circ\) has an odd Berge cycle. Then \(F\) is non-obligatory.

\end{proposition}
\begin{proof}
Suppress isolated vertices and put $m=|E(F)|$.  By Theorem~7.1 choose a graph
$G$ with $\chi(G)>\aleph_0$ and no odd cycle of length at most $m$.  Lemma~6.2
gives $\chi(\mathcal L(G))>\aleph_0$.

Suppose that $F$ embeds in $\mathcal L(G)$ and retain exactly the host
hyperedges corresponding to the edges of $F$.  The selected trace is finite
and linear.  Corollary~6.4 localises every Berge cycle in one expansion fibre
and sends a Berge cycle of length $\ell$ to an ordinary cycle of the same
length in $G$.  The given odd Berge cycle has $\ell\le m$, contradicting the
choice of $G$.  Thus $\mathcal L(G)$ is an uncountably chromatic $F$-free
host.
\end{proof}

The three propositions cover every failure of the intrinsic conditions, and prove

\[
\neg(3)\Longrightarrow\neg(1). \tag{8.2}
\]

\section{Proof of the classification}\label{conclusion}

Sections~3 and 4 prove \((2)\Longrightarrow(1)\). Section~5 proves
\((2)\Longleftrightarrow(3)\), and Section~8 proves the contrapositive
\(\neg(3)\Longrightarrow\neg(1)\). These implications establish all
equivalences in Theorem~A. In particular,

\begingroup\small
\[
\begin{gathered}
F\text{ is obligatory}
\iff
F\in\mathcal B,\\[1mm]
F\in\mathcal B
\iff
\begin{array}{l}
F^\circ\text{ is linear},\\[1mm]
\text{every hyperedge-node of }I(F^\circ)
\text{ is incident with a bridge},\\[1mm]
\text{every Berge cycle of }F^\circ\text{ has even length}.
\end{array}
\end{gathered}
\]
\endgroup

All embeddings are non-induced: additional host hyperedges on the image vertices are irrelevant throughout. The proof uses only ZFC, graph-colouring compactness, the classical Erd\H{o}s--Hajnal high-odd-girth theorem, and the standard Erd\H{o}s--Rado relation (8.1). It assumes neither CH nor GCH, and no forcing axiom or large-cardinal hypothesis.

\section{Finite parameter consequences}\label{finite-parameter-consequences}

In this section a triple system is \emph{connected} when its Levi graph is
connected.  We work first without isolated vertices.  For an integer $r\ge1$,
put
\[
q(r)=\left\lceil2\sqrt r\right\rceil.
\]

\begin{proposition}[Edge-deletion form of the bridge condition]
\label{proposition-edge-deletion-bridge-condition}
Let $F$ be a finite triple system and let $e=\{x,y,z\}\in E(F)$.  The
hyperedge-node $e$ has no incident bridge in $I(F)$ if and only if $x,y,z$
lie in one connected component after the hyperedge $e$ is deleted and all
points are retained.  Consequently the bridge condition in Theorem~A is
equivalent to the following hypergraph-native condition:
\[
\text{for every }e\in E(F^\circ),\text{ deleting }e\text{ separates its
three points into at least two components.}
\]
\end{proposition}
\begin{proof}
Fix $x\in e$.  The Levi incidence $xe$ is not a bridge precisely when there
is an alternative path from $e$ to $x$ after that incidence is deleted.  A
simple such path leaves $e$ through $y$ or $z$ and thereafter lies entirely
in the Levi graph of $F-e$.  Thus $xe$ is non-bridging exactly when $x$ is
connected in $F-e$ to at least one of the other two points of $e$.

If all three incidences at $e$ are non-bridging, no one of $x,y,z$ can form a
singleton component among those three points.  A partition of three points
with this property has one block, so all three are connected in $F-e$.  The
converse follows by reversing the same paths.
\end{proof}

\begin{lemma}[Bipartite shadow]
\label{lemma-bipartite-shadow}
Let $F$ be an obligatory finite triple system with no isolated vertices and
at least one hyperedge.  There is a finite bipartite graph $J$, without
isolated vertices, such that
\[
|E(J)|=|E(F)|,\qquad |V(J)|=|V(F)|-|E(F)|,
\]
and $J$ has the same number of connected components as $F$.  Conversely, for
every finite bipartite graph $J$ without isolated vertices, the expansion
$J^+$ is obligatory and has these parameter relations.
\end{lemma}
\begin{proof}
By Theorem~A and Proposition~\ref{proposition-5.2-bridge-block-decomposition},
each connected component $H$ of $F$ is obtained by one-point amalgamations
from connected expansion pieces
\[
J_1^+,\ldots,J_k^+,
\]
where each $J_i$ is a connected bipartite graph.  Disconnected graph pieces
may simply be split into their connected components.  Write
$m_i=|E(J_i)|$ and $s_i=|V(J_i)|$.  Since $H$ is connected, its $k$ pieces
are joined by exactly $k-1$ one-point identifications.  Hence
\[
|E(H)|=\sum_{i=1}^k m_i,
\qquad
|V(H)|=\sum_{i=1}^k(m_i+s_i)-(k-1).
\tag{10.1}
\]

Take one-point sums of the ordinary graphs $J_1,\ldots,J_k$ along any tree on
the $k$ indices.  Before each identification, interchange the two
bipartition classes of the new factor if necessary.  The resulting graph
$J_H$ is connected and bipartite, and (10.1) gives
\[
|E(J_H)|=|E(H)|,
\qquad
|V(J_H)|=|V(H)|-|E(H)|.
\]
Taking the disjoint union of these shadows over the components of $F$ proves
the forward assertion.  The converse is the expansion-generator case of
Theorem~A, with the displayed counts immediate from the definition of $J^+$.
\end{proof}

\begin{theorem}[Exact order--size--component spectrum]
\label{theorem-order-size-component-spectrum}
Let $m\ge1$, $1\le c\le m$, and $n$ be integers.  There exists an obligatory
triple system $F$ with no isolated vertices, exactly $m$ hyperedges, exactly
$n$ vertices, and exactly $c$ connected components if and only if
\[
\boxed{
 m+2(c-1)+\left\lceil2\sqrt{m-c+1}\right\rceil
 \le n\le 2m+c.
}
\tag{10.2}
\]
Every integer in this interval occurs.
\end{theorem}
\begin{proof}
We first record the corresponding elementary graph fact.  A connected simple
bipartite graph with $r\ge1$ edges and $s$ vertices exists if and only if
\[
q(r)\le s\le r+1.
\tag{10.3}
\]
Indeed, connectedness gives $r\ge s-1$.  If the bipartition has sizes $a,b$,
then
\[
r\le ab\le\left\lfloor\frac{s^2}{4}\right\rfloor,
\]
which is equivalent to $s\ge q(r)$.  Conversely,
$K_{\lfloor s/2\rfloor,\lceil s/2\rceil}$ contains a spanning tree, and
starting with that tree and adding arbitrary unused edges realises every edge
count between $s-1$ and $\lfloor s^2/4\rfloor$.

Apply Lemma~\ref{lemma-bipartite-shadow}.  Let the $c$ connected components of
the shadow have positive edge counts $m_1,\ldots,m_c$, where
$\sum_i m_i=m$.  If their orders are $s_i$, then (10.3) gives
\[
\sum_{i=1}^c q(m_i)\le \sum_{i=1}^c s_i\le m+c.
\tag{10.4}
\]
For positive integers $a,b$,
\[
q(a)+q(b)\ge q(a+b-1)+2.
\tag{10.5}
\]
To see this, note that
\[
\sqrt a+\sqrt b\ge\sqrt{a+b-1}+1,
\]
because, after squaring, the assertion reduces to
$(\sqrt a-1)(\sqrt b-1)\ge0$; taking ceilings gives (10.5).  Repeatedly
applying (10.5) to the left side of (10.4) yields
\[
\sum_{i=1}^c q(m_i)
\ge q(m-c+1)+2(c-1).
\]
The shadow has $n-m$ vertices, so these inequalities give (10.2).

For sharpness put $M=m-c+1$.  Take $c-1$ components equal to $K_2$.  By
(10.3), a final connected bipartite component with $M$ edges can have every
order from $q(M)$ through $M+1$.  Their disjoint union therefore realises
every shadow order from $q(M)+2(c-1)$ through $m+c$.  Expanding it adds
exactly $m$ private vertices and proves attainability of every integer in
(10.2).
\end{proof}

\begin{corollary}[Connected and unrestricted spectra]
\label{corollary-connected-order-size-spectrum}
For $m\ge1$:
\begin{enumerate}
\item a connected obligatory triple system without isolated vertices has
$m$ hyperedges and $n$ vertices if and only if
\[
m+\left\lceil2\sqrt m\right\rceil\le n\le2m+1;
\]
\item without prescribing the number of components, a reduced obligatory
triple system has $m$ hyperedges and $n$ vertices if and only if
\[
m+\left\lceil2\sqrt m\right\rceil\le n\le3m;
\]
\item if isolated vertices are permitted, an obligatory triple system with
$m$ hyperedges and $n$ vertices exists if and only if
\[
n\ge m+\left\lceil2\sqrt m\right\rceil.
\]
\end{enumerate}
\end{corollary}
\begin{proof}
The first assertion is Theorem~\ref{theorem-order-size-component-spectrum}
with $c=1$.  For the second, the connected interval covers through $2m+1$;
for $n>2m+1$, take $c=n-2m$ and use the upper endpoint $n=2m+c$ of (10.2).
The largest possible value is obtained from $m$ disjoint triples.  The final
assertion follows by adjoining arbitrary isolated vertices.
\end{proof}

\begin{corollary}[Size spectrum at fixed order]
\label{corollary-fixed-order-size-spectrum}
Let $c\ge1$ and $n\ge3c$.  A reduced obligatory triple system with exactly
$n$ vertices, $c$ connected components, and $m$ hyperedges exists if and only
if
\[
\boxed{
\left\lceil\frac{n-c}{2}\right\rceil
\le m\le
n-2c+4-\left\lceil2\sqrt{n-3c+4}\right\rceil.
}
\tag{10.6}
\]
Consequently, a reduced obligatory triple system with exactly $c$ connected
components and $n$ vertices exists if and only if
\[
n=3c,\qquad n=3c+2,\qquad\text{or}\qquad n\ge3c+4.
\tag{10.7}
\]
\end{corollary}
\begin{proof}
The upper-order inequality in (10.2) is equivalent to
$m\ge\lceil(n-c)/2\rceil$.  For the other inequality put
\[
M=m-c+1,\qquad N=n-3c+3.
\]
It becomes
\[
M+\left\lceil2\sqrt M\right\rceil\le N.
\tag{10.8}
\]
For integers $N\ge3$ and $M\ge1$, the largest $M$ satisfying (10.8) is
\[
U=N+2-\left\lceil2\sqrt{N+1}\right\rceil.
\tag{10.9}
\]
Indeed, write $k=\lceil2\sqrt{N+1}\rceil$.  Since
$k^2\ge4(N+1)$,
\[
4U=4N+8-4k\le(k-2)^2,
\]
so $\lceil2\sqrt U\rceil\le k-2$ and
$U+\lceil2\sqrt U\rceil\le N$.  On the other hand,
$(k-1)^2<4(N+1)$ gives
\[
(k-3)^2<4(N+3-k)=4(U+1),
\]
so $\lceil2\sqrt{U+1}\rceil\ge k-2$ and
$(U+1)+\lceil2\sqrt{U+1}\rceil\ge N+1$.  The left side of (10.8) is
strictly increasing in $M$, proving (10.9).  Translating back to $m$ gives
(10.6).

For (10.7), Corollary~\ref{corollary-connected-order-size-spectrum} shows that
a connected reduced obligatory system has order $3$, order $5$, or any order
at least $7$.  The first two statements come from $m=1,2$.  For $m\ge3$, the
connected intervals have no gaps because
$\lceil2\sqrt{m+1}\rceil\le m+1$ for $m+1\ge4$.  Thus every component has
order
\[
3+r,\qquad r\in\{0,2\}\cup[4,\infty).
\]
The sum of $c$ such increments is again $0$, $2$, or any integer at least
$4$: use one nonzero increment and take all remaining increments to be zero.
The increments $1$ and $3$ are impossible.  This proves (10.7).
\end{proof}

\begin{corollary}[Exact Levi cycle-rank spectrum]
\label{corollary-levi-cycle-rank-spectrum}
Let $F$ be reduced and obligatory, with $m$ hyperedges, $n$ vertices, and $c$
connected components.  The cyclomatic number of its Levi graph is
\[
\beta(I(F))=|E(I(F))|-|V(I(F))|+c=2m-n+c.
\tag{10.10}
\]
For fixed $m,c$, every integer in the interval
\[
0\le\beta(I(F))\le
m-c+2-\left\lceil2\sqrt{m-c+1}\right\rceil
\tag{10.11}
\]
occurs, and no other value occurs.  In particular,
\[
|V(F)|=2|E(F)|+c
\quad\Longleftrightarrow\quad
I(F)\text{ is a forest}.
\]
\end{corollary}
\begin{proof}
The Levi graph has $n+m$ vertices, $3m$ incidence edges, and $c$ components,
which gives (10.10).  Substitution of the two endpoints in (10.2) gives
(10.11), and every intermediate value occurs because every intermediate order
does.  A finite graph has cyclomatic number zero exactly when it is a forest.
\end{proof}

\begin{corollary}[Rigidity at balanced lower endpoints]
\label{corollary-balanced-endpoint-rigidity}
Let $F$ be connected, reduced, and obligatory.
\begin{enumerate}
\item If $|E(F)|=t^2$ and $|V(F)|=t^2+2t$, then
$F\cong K_{t,t}^+$.
\item If $|E(F)|=t(t+1)$ and $|V(F)|=t(t+1)+2t+1$, then
$F\cong K_{t,t+1}^+$.
\end{enumerate}
Here $t\ge1$, and the complete bipartite graphs are unique up to interchanging
the two classes.
\end{corollary}
\begin{proof}
In either case the shadow from Lemma~\ref{lemma-bipartite-shadow} attains
$|E(J)|=\lfloor|V(J)|^2/4\rfloor$.  Equality forces every possible cross edge
to be present and the two bipartition classes to be balanced, so the shadow is
$K_{t,t}$ or $K_{t,t+1}$, respectively.

For $t\ge2$ these complete bipartite graphs have no cut vertex.  On the other
hand, the shadow built from two or more connected expansion pieces is a
nontrivial one-point sum and has a cut vertex.  Hence the bridge-block
reconstruction has only one piece, and $F$ is the expansion of the displayed
complete bipartite graph.  The cases $t=1$ are immediate.
\end{proof}

\section*{Formal verification and reproducibility}

The self-contained Lean~4 source is available in the public
\href{https://github.com/SamPetkov/Erdos593/blob/main/formalization/Erdos593SelfContained.lean}{\texttt{Erdos593} repository}.
The public commit history records a checked formalisation scaffold on
\href{https://github.com/SamPetkov/Erdos593/commit/ba37b8c511ab390c14a905110366d7cac1ffa08f}{15 July 2026}
and the complete finite structural classification later that day in
\href{https://github.com/SamPetkov/Erdos593/commit/6fd00a76064401f3f10aabef474f59d3c6ecd6bf}{commit \texttt{6fd00a7}}.
To the author's knowledge, this is the first publicly timestamped Lean
formalisation of the finite structural theorem.  Li's arXiv v2, submitted on
23 July 2026, contains a separate formal verification of the broader results
in his paper.  No assertion of shared code or derivation between the two Lean
developments is made.
For every finite triple system $F$, the exported theorems establish
\[
\begin{aligned}
\mathtt{F.IsObligatory}
&\Longleftrightarrow \mathtt{F.isolatedReduction.Intrinsic},\\
\mathtt{F.IsObligatory}
&\Longleftrightarrow \mathtt{Constructible\ F.isolatedReduction}.
\end{aligned}
\]
Thus the development verifies both directions of the classification,
including the isolated-vertex reduction.  Here ``finite'' refers to the
system $F$ being classified; host triple systems are not assumed finite and
are quantified within the formalisation's documented ambient-universe
convention.  The formalisation supplements, rather than replaces,
mathematical review.

\section*{Acknowledgments}

The author thanks Tom de Groot for his advice on revising the manuscript for
greater clarity.  The author also thanks Eric Li for a discussion of the
relationship between the two proofs.

\section*{AI assistance}

OpenAI's GPT-5.6 Pro through ChatGPT and Aristotle~\citep{achim-et-al-2025}
were used for proof development, adversarial checking, editorial restructuring,
and the Lean formalisation.  The author reviewed all incorporated suggestions
and assumes full responsibility for the arguments, citations, and final
manuscript.

\section*{Funding}
The author received no funding for this work.

\section*{Competing interests}
The author declares no competing interests.

\begin{thebibliography}{13}
\providecommand{\natexlab}[1]{#1}
\providecommand{\url}[1]{\texttt{#1}}
\expandafter\ifx\csname urlstyle\endcsname\relax
  \providecommand{\doi}[1]{doi: #1}\else
  \providecommand{\doi}{doi: \begingroup \urlstyle{rm}\Url}\fi

\bibitem[Erd{\H{o}}s(1995)]{erdos1995}
Paul Erd{\H{o}}s.
\newblock On some problems in combinatorial set theory.
\newblock \emph{Publikacije Instituta Matemati\v{c}kog (Beograd) (N.S.)},
  57(71):\penalty0 61--65, 1995.

\bibitem[Bloom(2026)]{bloom593}
Thomas~F. Bloom.
\newblock Erd{\H{o}}s problem 593.
\newblock \url{https://www.erdosproblems.com/593}, 2026.
\newblock Accessed 11 July 2026.

\bibitem[Li(2026)]{li2026}
Eric Li.
\newblock A resolution of {Erd{\H{o}}s Problems 593 and 1177}: Obligatory
  triple systems and exact spectra.
\newblock arXiv:2606.24882, submitted 23 June 2026.
\newblock \doi{10.48550/arXiv.2606.24882}.
\newblock URL \url{https://arxiv.org/abs/2606.24882}.

\bibitem[Erd{\H{o}}s and Hajnal(1966)]{erdoshajnal1966}
Paul Erd{\H{o}}s and Andr{\'a}s Hajnal.
\newblock On chromatic number of graphs and set-systems.
\newblock \emph{Acta Mathematica Academiae Scientiarum Hungaricae}, 17\penalty0
  (1--2):\penalty0 61--99, 1966.
\newblock \doi{10.1007/BF02020444}.

\bibitem[Erd{\H{o}}s et~al.(1973)Erd{\H{o}}s, Hajnal, and Rothschild]{ehr1973}
Paul Erd{\H{o}}s, Andr{\'a}s Hajnal, and Bruce~L. Rothschild.
\newblock On chromatic number of graphs and set-systems.
\newblock In \emph{Cambridge Summer School in Mathematical Logic (Cambridge,
  1971)}, volume 337 of \emph{Lecture Notes in Mathematics}, pages 531--538.
  Springer, Berlin and New York, 1973.

\bibitem[Erd{\H{o}}s et~al.(1975)Erd{\H{o}}s, Galvin, and Hajnal]{egh1975}
Paul Erd{\H{o}}s, Fred Galvin, and Andr{\'a}s Hajnal.
\newblock On set-systems having large chromatic number and not containing
  prescribed subsystems.
\newblock In \emph{Infinite and Finite Sets (Keszthely, 1973)}, volume~10 of
  \emph{Colloquia Mathematica Societatis J{\'a}nos Bolyai}, pages 425--513.
  North-Holland, 1975.

\bibitem[Komj{\'a}th(2001)]{komjath2001}
P{\'e}ter Komj{\'a}th.
\newblock Some remarks on obligatory subsystems of uncountably chromatic triple
  systems.
\newblock \emph{Combinatorica}, 21\penalty0 (2):\penalty0 233--238, 2001.
\newblock \doi{10.1007/s004930100021}.

\bibitem[Hajnal and Komj{\'a}th(2008)]{hajnalkomjath2008}
Andr{\'a}s Hajnal and P{\'e}ter Komj{\'a}th.
\newblock Obligatory subsystems of triple systems.
\newblock \emph{Acta Mathematica Hungarica}, 119\penalty0 (1--2):\penalty0
  1--13, 2008.
\newblock \doi{10.1007/s10474-007-6231-2}.

\bibitem[Komj{\'a}th(2008)]{komjath2008}
P{\'e}ter Komj{\'a}th.
\newblock An uncountably chromatic triple system.
\newblock \emph{Acta Mathematica Hungarica}, 121\penalty0 (1--2):\penalty0
  79--92, 2008.
\newblock \doi{10.1007/s10474-008-7179-6}.

\bibitem[Reiher(2024)]{reiher2024}
Christian Reiher.
\newblock Obligatory hypergraphs.
\newblock \emph{Proceedings of the American Mathematical Society}, 2024.
\newblock \doi{10.1090/proc/17021}.
\newblock URL \url{https://arxiv.org/abs/2403.11223}.
\newblock arXiv:2403.11223.

\bibitem[de~Bruijn and Erd{\H{o}}s(1951)]{debruijn1951}
Nicolaas~Govert de~Bruijn and Paul Erd{\H{o}}s.
\newblock A colour problem for infinite graphs and a problem in the theory of
  relations.
\newblock \emph{Nederl. Akad. Wetensch. Proc. Ser. A}, 54\penalty0
  (5):\penalty0 371--373, 1951.
\newblock \doi{10.1016/S1385-7258(51)50053-7}.

\bibitem[Erd{\H{o}}s and Rado(1956)]{erdosrado1956}
Paul Erd{\H{o}}s and Richard Rado.
\newblock A partition calculus in set theory.
\newblock \emph{Bulletin of the American Mathematical Society}, 62\penalty0
  (5):\penalty0 427--489, 1956.
\newblock \doi{10.1090/S0002-9904-1956-10036-0}.

\bibitem[Achim et~al.(2025)Achim, Best, Der, F{\'e}d{\'e}rico, Gukov,
  Halpern-Leistner, Henningsgard, Kudryashov, Meiburg, Michelsen, Patterson,
  Rodriguez, Scharff, Shanker, Sicca, Sowrirajan, Swope, Tamas, Tenev, Thomm,
  Williams, and Wu]{achim-et-al-2025}
Tudor Achim, Alex Best, Kevin Der, Math{\"i}s F{\'e}d{\'e}rico, Sergei Gukov,
  Daniel Halpern-Leistner, Kirsten Henningsgard, Yury Kudryashov, Alexander
  Meiburg, Martin Michelsen, Riley Patterson, Eric Rodriguez, Laura Scharff,
  Vikram Shanker, Vladimir Sicca, Hari Sowrirajan, Aidan Swope, Matyas Tamas,
  Vlad Tenev, Jonathan Thomm, Harold Williams, and Lawrence Wu.
\newblock {Aristotle}: {IMO}-level automated theorem proving, 2025.
\newblock URL \url{https://arxiv.org/abs/2510.01346}.
\newblock arXiv:2510.01346.

\end{thebibliography}

\end{document}
